\chapter{Kitaev Quantum Double Models}
\label{chap:kitaev-quantum-double-models}

The toric code is the quantum double model for the finite group
\(G=\ZZ_2\).  The purpose of this chapter is to explain the general finite-group construction.  The resulting models, usually denoted by \(D(G)\), are
commuting-projector Hamiltonians whose local Hilbert spaces are group algebras
rather than qubits.  They are exactly solvable for every finite group \(G\), and
for non-Abelian \(G\) they give fixed-point lattice models with non-Abelian
anyon excitations.

There are two guiding ideas.
\begin{enumerate}[label=(\roman*)]
    \item The toric-code star operator is a local \(\ZZ_2\) gauge
    transformation.  For a general finite group \(G\), the vertex term is the
    projector onto local gauge-invariant states.
    \item The toric-code plaquette operator enforces zero \(\ZZ_2\) flux.  For a
    general finite group \(G\), the plaquette term projects onto states with
    trivial group holonomy around each plaquette.
\end{enumerate}
The exact solvability follows from the same mechanism as in the toric code:
local gauge projectors and flatness projectors commute.

\section{From \texorpdfstring{\(\ZZ_2\)}{Z2} to finite groups}
\label{sec:quantum-double-from-z2-to-g}

Let \(G\) be a finite group with identity element \(e\).  The local Hilbert
space on each oriented edge \(\ell\) is the group algebra
\begin{equation}
\label{eq:edge-hilbert-group-algebra}
    \calH_{\ell}=\CC[G]
    =\operatorname{span}\{\ket{g}:g\in G\},
    \qquad
    \braket{g}{g'}=\delta_{g,g'} .
\end{equation}
The total Hilbert space on a lattice \(\Gamma=(V,E,P)\), with vertices \(V\),
edges \(E\), and plaquettes \(P\), is
\begin{equation}
\label{eq:quantum-double-total-hilbert-space}
    \calH_G=\bigotimes_{\ell\in E}\CC[G]_{\ell},
    \qquad
    \dim \calH_G=\abs{G}^{\abs{E}} .
\end{equation}
A basis state is a group-valued lattice gauge field
\begin{equation}
\label{eq:group-lattice-configuration}
    \ket{\{g_{\ell}\}}
    =\bigotimes_{\ell\in E}\ket{g_{\ell}}_{\ell} .
\end{equation}
Each edge is assigned an arbitrary reference orientation.  If a path traverses
an edge \(\ell\) in the reference direction, it contributes \(g_{\ell}\); if it
traverses \(\ell\) in the opposite direction, it contributes \(g_{\ell}^{-1}\).
This is the lattice analogue of parallel transport by a gauge connection.

For every \(h\in G\), define left and right multiplication operators on
\(\CC[G]\) by
\begin{equation}
\label{eq:left-right-regular-actions}
    L^h\ket{g}=\ket{hg},
    \qquad
    R^h\ket{g}=\ket{g h^{-1}} .
\end{equation}
The inverse in the definition of \(R^h\) is convenient because both \(L^h\) and
\(R^h\) define honest group actions:
\begin{equation}
\label{eq:left-right-group-actions}
    L^hL^k=L^{hk},
    \qquad
    R^hR^k=R^{hk},
    \qquad
    [L^h,R^k]=0 .
\end{equation}

\begin{remark}[Orientation convention]
A local gauge transformation at a vertex \(v\) should multiply an outgoing edge
variable from the left and an incoming edge variable from the right.  With the
convention \cref{eq:left-right-regular-actions}, this means
\begin{equation}
\label{eq:gauge-transform-edge-variable}
    g_{\ell}\longmapsto
    \begin{cases}
        h g_{\ell}, & \ell\text{ starts at }v,\\
        g_{\ell}h^{-1}, & \ell\text{ ends at }v.
    \end{cases}
\end{equation}
This is the lattice version of \(U_{ij}\mapsto h_iU_{ij}h_j^{-1}\).
\end{remark}

\section{Vertex gauge transformations}
\label{sec:quantum-double-vertex-operators}

For a vertex \(v\in V\) and a group element \(h\in G\), define
\begin{equation}
\label{eq:vertex-gauge-transformation}
    A_v^h
    =
    \prod_{\ell:\,s(\ell)=v} L_{\ell}^{h}
    \prod_{\ell:\,t(\ell)=v} R_{\ell}^{h},
\end{equation}
where \(s(\ell)\) and \(t(\ell)\) denote the source and target of the oriented
edge \(\ell\).  Thus \(A_v^h\) implements a gauge transformation by \(h\) at the
single vertex \(v\).  The vertex projector is the group average
\begin{equation}
\label{eq:quantum-double-vertex-projector}
    A_v=\frac{1}{\abs{G}}\sum_{h\in G}A_v^h .
\end{equation}

\begin{proposition}[Vertex projectors]
The operators \(A_v\) are mutually commuting projectors:
\begin{equation}
\label{eq:vertex-projector-algebra}
    A_v^2=A_v,
    \qquad
    [A_v,A_w]=0 .
\end{equation}
\end{proposition}

\begin{proof}
For a fixed vertex \(v\), the map \(h\mapsto A_v^h\) is a representation of
\(G\), so
\begin{equation}
\label{eq:vertex-projector-proof-1}
    A_v^2
    =\frac{1}{\abs{G}^2}\sum_{h,k\in G}A_v^{hk}
    =\frac{1}{\abs{G}}\sum_{q\in G}A_v^q
    =A_v .
\end{equation}
For \(v\neq w\), the only possible overlap is an edge connecting \(v\) and
\(w\).  On such an edge one operator acts by left multiplication and the other
by right multiplication.  These commute by \cref{eq:left-right-group-actions}.
Therefore \([A_v^h,A_w^k]=0\) for all \(h,k\), and hence
\([A_v,A_w]=0\).
\end{proof}

The condition
\begin{equation}
\label{eq:gauge-invariant-state-condition}
    A_v\ket{\psi}=\ket{\psi}
\end{equation}
is the lattice Gauss-law constraint.  Thus the vertex terms impose local gauge
invariance.

\section{Plaquette holonomy and flatness projectors}
\label{sec:quantum-double-plaquette-operators}

Choose an orientation and a base vertex \(v\in\partial p\) for each plaquette
\(p\).  Let \(\partial p=(\ell_1,\ldots,\ell_n)\) be the ordered boundary path
starting and ending at \(v\).  Define the plaquette holonomy
\begin{equation}
\label{eq:plaquette-holonomy}
    H_{p,v}(\{g_{\ell}\})
    =
    g_{\ell_1}^{\epsilon_1}
    g_{\ell_2}^{\epsilon_2}
    \cdots
    g_{\ell_n}^{\epsilon_n},
    \qquad
    \epsilon_j=\begin{cases}
        +1, & \ell_j\text{ is traversed with the reference orientation},\\
        -1, & \ell_j\text{ is traversed against it}.
    \end{cases}
\end{equation}
For \(k\in G\), define a plaquette flux projector by
\begin{equation}
\label{eq:plaquette-flux-projector-k}
    B_p^k\ket{\{g_{\ell}\}}
    =
    \delta_{H_{p,v}(\{g_{\ell}\}),k}\ket{\{g_{\ell}\}} .
\end{equation}
The Hamiltonian uses the trivial-flux projector
\begin{equation}
\label{eq:quantum-double-plaquette-projector}
    B_p:=B_p^e .
\end{equation}

\begin{remark}[Basepoint dependence]
For non-Abelian \(G\), the holonomy around a plaquette depends on the basepoint
by conjugation.  If the starting point is changed, \(H_{p,v}\) is replaced by
\(xH_{p,v}x^{-1}\) for some boundary transport \(x\).  Therefore the condition
\(H_{p,v}=e\) is basepoint-independent, while the condition \(H_{p,v}=k\) is not
basepoint-independent unless it is stated up to conjugacy.
\end{remark}

The plaquette projectors are diagonal in the group basis and obey
\begin{equation}
\label{eq:plaquette-projector-algebra}
    (B_p^k)^2=B_p^k,
    \qquad
    B_p^kB_p^{k'}=\delta_{k,k'}B_p^k,
    \qquad
    \sum_{k\in G}B_p^k=\id .
\end{equation}
The physical plaquette term \(B_p=B_p^e\) enforces flatness of the lattice gauge
field at \(p\).

\section{The quantum double Hamiltonian}
\label{sec:quantum-double-hamiltonian}

The Kitaev quantum double Hamiltonian is
\begin{equation}
\label{eq:quantum-double-hamiltonian}
    H_{D(G)}
    =
    -J_e\sum_{v\in V}A_v
    -J_m\sum_{p\in P}B_p,
    \qquad
    J_e,J_m>0 .
\end{equation}
Equivalently, up to an additive constant,
\begin{equation}
\label{eq:quantum-double-projector-form}
    H_{D(G)}
    =
    J_e\sum_{v\in V}(\id-A_v)
    +J_m\sum_{p\in P}(\id-B_p)
    -(J_e\abs{V}+J_m\abs{P}) .
\end{equation}
The exact ground states are those satisfying
\begin{equation}
\label{eq:quantum-double-ground-state-conditions}
    A_v\ket{\psi}=\ket{\psi},
    \qquad
    B_p\ket{\psi}=\ket{\psi},
    \qquad
    \forall v\in V,\ p\in P .
\end{equation}
Thus the ground space consists of gauge-invariant wavefunctions supported on
flat \(G\)-connections.

\begin{theorem}[Commuting-projector structure]
The terms of \(H_{D(G)}\) commute:
\begin{equation}
\label{eq:quantum-double-commuting-projectors}
    [A_v,A_w]=0,
    \qquad
    [B_p,B_q]=0,
    \qquad
    [A_v,B_p]=0 .
\end{equation}
Therefore \(H_{D(G)}\) is an exactly solvable commuting-projector Hamiltonian.
\end{theorem}

\begin{proof}
The first equality was proved in \cref{sec:quantum-double-vertex-operators}.  The
second equality follows because all \(B_p\) are diagonal multiplication
operators in the group basis.

It remains to check \([A_v,B_p]=0\).  If \(v\notin\partial p\), the operators act
on disjoint local data or act in a way that does not change the plaquette
holonomy, so they commute.  If \(v\in\partial p\), then a gauge transformation by
\(h\) at \(v\) conjugates the plaquette holonomy based at \(v\):
\begin{equation}
\label{eq:holonomy-conjugation-under-gauge-transform}
    H_{p,v}\longmapsto hH_{p,v}h^{-1} .
\end{equation}
The property \(H_{p,v}=e\) is invariant under conjugation.  Hence
\(A_v^hB_p^e=B_p^eA_v^h\) for every \(h\), and averaging over \(h\) gives
\([A_v,B_p]=0\).
\end{proof}

\begin{remark}[Why the identity flux is special]
The operator \(B_p^k\) with fixed noncentral \(k\) does not commute with a gauge
transformation at the basepoint.  Instead,
\begin{equation}
\label{eq:flux-projector-conjugation}
    A_v^h B_p^k A_v^{h^{-1}}=B_p^{hkh^{-1}} .
\end{equation}
The gauge-invariant flux observable is therefore the conjugacy-class projector
\begin{equation}
\label{eq:conjugacy-class-flux-projector}
    B_p^C=\sum_{k\in C}B_p^k,
\end{equation}
where \(C\subset G\) is a conjugacy class.  The Hamiltonian chooses the trivial
class \(C=\{e\}\).
\end{remark}

\section{Recovery of the toric code}
\label{sec:quantum-double-recovers-toric-code}

Set \(G=\ZZ_2=\{0,1\}\) with addition modulo two.  The group algebra
\(\CC[\ZZ_2]\) is a qubit Hilbert space with computational basis
\(\ket{0},\ket{1}\).  Left multiplication by the nontrivial group element flips
the basis state:
\begin{equation}
\label{eq:z2-left-multiplication-sigma-x}
    L^1\ket{g}=\ket{g+1}=\sigma^x\ket{g} .
\end{equation}
The vertex projector becomes
\begin{equation}
\label{eq:z2-vertex-projector-toric-code}
    A_v
    =\frac{1}{2}\left(\id+\prod_{\ell\ni v}\sigma^x_{\ell}\right).
\end{equation}
The plaquette flatness condition says that the sum of edge variables around a
plaquette is zero modulo two.  In the basis where
\(\sigma^z\ket{g}=(-1)^g\ket{g}\), this is
\begin{equation}
\label{eq:z2-plaquette-projector-toric-code}
    B_p
    =\frac{1}{2}\left(\id+\prod_{\ell\in\partial p}\sigma^z_{\ell}\right).
\end{equation}
Therefore \cref{eq:quantum-double-hamiltonian} becomes the toric-code
Hamiltonian up to an additive constant and an overall normalization of the
couplings.  In this precise sense,
\begin{equation}
\label{eq:toric-code-as-dz2}
    \text{toric code}=D(\ZZ_2) .
\end{equation}

\section{Ground states as flat gauge fields}
\label{sec:quantum-double-ground-states-flat-connections}

Let the lattice be embedded on a closed orientable surface \(\Sigma_g\).  The
constraint \(B_p=1\) imposes trivial holonomy around every contractible
plaquette.  Therefore a ground-state configuration is a flat lattice
\(G\)-connection.  The constraint \(A_v=1\) then identifies configurations
related by local gauge transformations.

Equivalently, the ground-state data are representations of the fundamental
group into \(G\), modulo global conjugation:
\begin{equation}
\label{eq:flat-g-connections-moduli}
    \operatorname{Flat}_G(\Sigma_g)
    \simeq
    \operatorname{Hom}(\pi_1(\Sigma_g),G)/G .
\end{equation}
Here \(G\) acts by simultaneous conjugation.  The ground-state Hilbert space is
the space of gauge-invariant complex functions on flat connections:
\begin{equation}
\label{eq:quantum-double-ground-space-functions}
    \calH_{\rm GS}(\Sigma_g)
    \simeq
    \operatorname{Fun}\!\bigl(\operatorname{Hom}(\pi_1(\Sigma_g),G)\bigr)^G .
\end{equation}
Thus
\begin{equation}
\label{eq:quantum-double-ground-degeneracy-orbits}
    \dim \calH_{\rm GS}(\Sigma_g)
    =
    \#\bigl(\operatorname{Hom}(\pi_1(\Sigma_g),G)/G\bigr),
\end{equation}
where the right-hand side denotes the number of conjugation orbits of flat
\(G\)-connections.

For a torus \(T^2\),
\begin{equation}
\label{eq:torus-fundamental-group-commuting-pairs}
    \pi_1(T^2)=\langle a,b:ab=ba\rangle .
\end{equation}
Therefore a flat connection is a commuting pair \((x,y)\in G\times G\), modulo
simultaneous conjugation:
\begin{equation}
\label{eq:torus-ground-state-commuting-pairs}
    \dim\calH_{\rm GS}(T^2)
    =
    \#\{(x,y)\in G\times G:xy=yx\}/G .
\end{equation}
For Abelian \(G\), conjugation is trivial, and this reduces to
\begin{equation}
\label{eq:abelian-quantum-double-ground-degeneracy}
    \dim\calH_{\rm GS}(\Sigma_g)=\abs{G}^{2g} .
\end{equation}
Taking \(G=\ZZ_2\) gives the toric-code result \(2^{2g}\).

\begin{remark}[Ground-state counting is not stabilizer counting]
For \(G=\ZZ_2\), the model is a Pauli stabilizer code and the degeneracy can be
counted by independent stabilizers.  For a general finite non-Abelian group,
there is no Pauli stabilizer structure.  The correct replacement is the counting
of gauge-invariant flat connections, or equivalently the representation theory
of the quantum double algebra.
\end{remark}

\section{Electric, magnetic, and dyonic excitations}
\label{sec:quantum-double-excitations}

Excitations are local violations of the commuting projectors.
\begin{itemize}[leftmargin=2.2em]
    \item A violation of \(A_v\) is an electric charge.  It transforms
    nontrivially under local gauge transformations.
    \item A violation of \(B_p\) is a magnetic flux.  Its gauge-invariant label
    is a conjugacy class of plaquette holonomy.
    \item In a non-Abelian group, charge and flux cannot in general be separated
    into independent labels.  The elementary superselection sectors are dyons.
\end{itemize}
The precise label of a dyon is a pair
\begin{equation}
\label{eq:quantum-double-anyon-label-pair}
    (C,\rho),
\end{equation}
where \(C\subset G\) is a conjugacy class and \(\rho\) is an irreducible
representation of the centralizer of a representative element of \(C\).

To define this label, choose an element \(c\in C\).  Its centralizer is
\begin{equation}
\label{eq:centralizer-definition}
    Z_c
    =
    \{h\in G:hc=ch\} .
\end{equation}
If \(c' = xcx^{-1}\) is another representative of the same conjugacy class, then
\(Z_{c'}=xZ_cx^{-1}\), so the set of irreducible representations is naturally
identified up to conjugation.  Hence the anyon label is well-defined as
\begin{equation}
\label{eq:anyon-label-set-dg}
    \mathcal{A}_{D(G)}
    =
    \{(C,\rho): C\in\operatorname{Conj}(G),\
    \rho\in\operatorname{Irr}(Z_c),\ c\in C\} .
\end{equation}

\begin{definition}[Quantum-double anyon]
An anyon type of the quantum double model \(D(G)\) is a pair \((C,\rho)\), where
\(C\) is a magnetic flux conjugacy class and \(\rho\) is an electric charge
representation of the centralizer of that flux.
\end{definition}

This formula has several important special cases.
\begin{enumerate}[label=(\roman*)]
    \item The vacuum is \((\{e\},\mathbf{1})\), where \(\mathbf{1}\) is the
    trivial representation of \(G\).
    \item Pure electric charges have trivial flux:
    \begin{equation}
    \label{eq:pure-electric-labels}
        (\{e\},\rho),
        \qquad
        \rho\in\operatorname{Irr}(G).
    \end{equation}
    \item Pure magnetic fluxes have nontrivial conjugacy class but trivial
    centralizer representation:
    \begin{equation}
    \label{eq:pure-magnetic-labels}
        (C,\mathbf{1}_{Z_c}) .
    \end{equation}
    \item Genuine dyons have both nontrivial flux and nontrivial charge
    representation.
\end{enumerate}

\section{The local quantum double algebra}
\label{sec:local-quantum-double-algebra}

The origin of the label \((C,\rho)\) is the finite-dimensional Hopf algebra known
as the Drinfeld double \(D(G)\).  At a puncture or excitation site, there are two
kinds of local operations: measuring flux and applying gauge transformations.
They obey a crossed relation.

Let \(P^g\) denote the projector onto flux \(g\), and let \(T^h\) denote a gauge
transformation by \(h\).  Then
\begin{equation}
\label{eq:local-double-crossed-relation}
    T^h P^g = P^{hgh^{-1}}T^h .
\end{equation}
Thus flux labels are not invariant under gauge transformations; only conjugacy
classes are invariant.  The algebra generated by \(P^g\) and \(T^h\), with
\begin{equation}
\label{eq:local-double-algebra-relations}
    P^gP^{g'}=\delta_{g,g'}P^g,
    \qquad
    T^hT^{h'}=T^{hh'},
    \qquad
    \sum_{g\in G}P^g=\id,
\end{equation}
is the quantum double algebra \(D(G)\) in its standard lattice form.

\begin{proposition}[Irreducible sectors of \(D(G)\)]
The irreducible representations of the local quantum double algebra are labelled
by pairs \((C,\rho)\), where \(C\) is a conjugacy class of \(G\), and \(\rho\) is
an irreducible representation of the centralizer \(Z_c\) of any representative
\(c\in C\).
\end{proposition}

\begin{proof}[Sketch]
First diagonalize the commuting flux projectors \(P^g\).  The gauge action
\(T^h\) moves a flux \(g\) to the conjugate flux \(hgh^{-1}\).  Therefore an
irreducible representation cannot be supported on a single noncentral group
element; it is supported on an entire conjugacy class \(C\).  After choosing a
representative \(c\in C\), the transformations that preserve this representative
are precisely the centralizer elements \(Z_c\).  The remaining internal charge
quantum numbers must therefore form an irreducible representation \(\rho\) of
\(Z_c\).  This gives the pair \((C,\rho)\).
\end{proof}

The quantum dimension of the anyon \((C,\rho)\) is
\begin{equation}
\label{eq:quantum-double-quantum-dimension}
    d_{(C,\rho)}=\abs{C}\,\dim\rho .
\end{equation}
The total quantum dimension is
\begin{equation}
\label{eq:quantum-double-total-quantum-dimension}
    \mathcal{D}
    =
    \sqrt{\sum_{(C,\rho)}d_{(C,\rho)}^2}
    =\abs{G} .
\end{equation}
Indeed,
\begin{equation}
\label{eq:quantum-double-total-dimension-check}
    \sum_{(C,\rho)}d_{(C,\rho)}^2
    =
    \sum_C \abs{C}^2\sum_{\rho\in\operatorname{Irr}(Z_c)}(\dim\rho)^2
    =
    \sum_C \abs{C}^2\abs{Z_c}
    =
    \sum_C \abs{C}\abs{G}
    =
    \abs{G}^2 .
\end{equation}

The topological spin of \((C,\rho)\) is
\begin{equation}
\label{eq:quantum-double-topological-spin}
    \theta_{(C,\rho)}
    =
    \frac{\chi_{\rho}(c)}{\dim\rho},
    \qquad c\in C,
\end{equation}
where \(\chi_{\rho}\) is the character of \(\rho\).  Since \(c\) lies in the
center of its own centralizer \(Z_c\), Schur's lemma ensures that this ratio is a
phase for a unitary irreducible representation.

\section{Abelian case and the toric-code anyons}
\label{sec:abelian-quantum-double-anyons}

If \(G\) is Abelian, every conjugacy class contains one element and every
centralizer is the whole group.  Therefore
\begin{equation}
\label{eq:abelian-anyon-labels}
    \mathcal{A}_{D(G)}
    =
    G\times \widehat{G},
\end{equation}
where \(\widehat{G}=\operatorname{Irr}(G)\) is the character group.  An anyon is
labelled by a magnetic flux \(g\in G\) and an electric character
\(\chi\in\widehat{G}\).  All anyons have quantum dimension one:
\begin{equation}
\label{eq:abelian-quantum-double-dimensions}
    d_{(g,\chi)}=1,
    \qquad
    \mathcal{D}=\abs{G} .
\end{equation}
The braiding phase between \((g,\chi)\) and \((h,\psi)\) is determined by the
Aharonov--Bohm pairing between charge and flux:
\begin{equation}
\label{eq:abelian-double-mutual-braiding}
    M_{(g,\chi),(h,\psi)}
    =
    \chi(h)\,\psi(g) .
\end{equation}

For \(G=\ZZ_2\), there are two group elements \(0,1\) and two characters
\(\mathbf{1},\chi\), with \(\chi(1)=-1\).  The four anyons are
\begin{equation}
\label{eq:toric-code-anyon-labels-from-dz2}
    1=(0,\mathbf{1}),
    \qquad
    e=(0,\chi),
    \qquad
    m=(1,\mathbf{1}),
    \qquad
    \epsilon=(1,\chi).
\end{equation}
This reproduces the toric-code fusion rules
\begin{equation}
\label{eq:toric-code-fusion-from-dz2}
    e^2=m^2=\epsilon^2=1,
    \qquad
    e\times m=\epsilon .
\end{equation}
The mutual phase between \(e\) and \(m\) is \(-1\), because
\begin{equation}
\label{eq:toric-code-mutual-phase-from-character}
    M_{e,m}=\chi(1)=-1 .
\end{equation}

\section{Ribbon operators}
\label{sec:quantum-double-ribbon-operators}

In the toric code, open string operators create pairs of anyons at their
endpoints.  In the non-Abelian quantum double model, the correct generalization
is a ribbon operator.  A ribbon is a strip made from a direct-lattice path and a
nearby dual-lattice path.  It carries both electric and magnetic information.

A ribbon operator has two conceptual components.
\begin{enumerate}[label=(\roman*)]
    \item A magnetic component changes group variables on crossed edges and
    creates plaquette fluxes at the ribbon endpoints.
    \item An electric component parallel-transports representation indices along
    the direct path and creates gauge-charge violations at the endpoints.
\end{enumerate}
Closed ribbon operators commute with the Hamiltonian and measure topological
charge.  Open ribbon operators commute with all projectors away from their
endpoints and create a pair of excitations whose total topological charge is
trivial.

For pure electric probes, one obtains Wilson-loop operators.  If \(\gamma\) is a
closed oriented loop and \(R\) is a representation of \(G\), define
\begin{equation}
\label{eq:quantum-double-wilson-loop}
    W_R(\gamma)
    =
    \chi_R\!\left(\prod_{\ell\in\gamma}g_{\ell}^{\epsilon_{\ell}}\right),
\end{equation}
where \(\chi_R\) is the character of \(R\).  The trace makes
\(W_R(\gamma)\) invariant under conjugation of the loop holonomy.  On a
contractible loop in the ground state, flatness forces \(W_R(\gamma)=\dim R\).
On noncontractible loops, Wilson operators distinguish different topological
sectors.

\begin{remark}[Why ribbons replace strings]
For Abelian \(G\), electric and magnetic string operators can be separated
cleanly, and their algebra is essentially a phase.  For non-Abelian \(G\), moving
flux changes the charge frame by conjugation.  A ribbon keeps track of both the
flux transport and the charge indices.  This is why the excitation label is not
``flux plus representation of \(G\)'', but rather ``flux conjugacy class plus
representation of the corresponding centralizer.''
\end{remark}

\section{Example: \texorpdfstring{\(D(S_3)\)}{D(S3)}}
\label{sec:quantum-double-s3-example}

The smallest non-Abelian example is \(G=S_3\), the permutation group of three
objects.  Its conjugacy classes are
\begin{equation}
\label{eq:s3-conjugacy-classes}
    C_e=\{e\},
    \qquad
    C_t=\{(12),(23),(13)\},
    \qquad
    C_c=\{(123),(132)\} .
\end{equation}
The corresponding centralizers have sizes
\begin{equation}
\label{eq:s3-centralizer-sizes}
    \abs{Z_e}=6,
    \qquad
    \abs{Z_{(12)}}=2,
    \qquad
    \abs{Z_{(123)}}=3 .
\end{equation}
The irreducible representations are therefore distributed as follows:
\begin{equation}
\label{eq:s3-anyon-count}
    \#\mathcal{A}_{D(S_3)}
    =
    \#\operatorname{Irr}(S_3)
    +\#\operatorname{Irr}(\ZZ_2)
    +\#\operatorname{Irr}(\ZZ_3)
    =3+2+3=8 .
\end{equation}
Thus \(D(S_3)\) has eight anyon types.  Their quantum dimensions are
\begin{equation}
\label{eq:s3-quantum-dimensions-list}
    1,1,2,
    \qquad
    3,3,
    \qquad
    2,2,2,
\end{equation}
where the first group comes from pure charges over \(C_e\), the second from
transposition fluxes, and the third from three-cycle fluxes.  The dimension
check gives
\begin{equation}
\label{eq:s3-total-quantum-dimension-check}
    1^2+1^2+2^2+3^2+3^2+2^2+2^2+2^2=36=\abs{S_3}^2 .
\end{equation}
This example illustrates the genuinely non-Abelian feature: some anyons have
quantum dimension larger than one.

\section{Relation to topological quantum field theory}
\label{sec:quantum-double-tqft-viewpoint}

The model \(D(G)\) is the Hamiltonian lattice realization of finite-group
Dijkgraaf--Witten gauge theory with trivial cocycle.  Its input data are purely
group-theoretic, but its low-energy physics is topological.  The ground states
only depend on the topology of the spatial surface, and quasiparticles are
classified by the representation category of the Drinfeld double,
\begin{equation}
\label{eq:quantum-double-category}
    \text{anyon category of }D(G)
    =
    \operatorname{Rep}(D(G)) .
\end{equation}
The toric code corresponds to \(\operatorname{Rep}(D(\ZZ_2))\).

The commuting-projector Hamiltonian should be viewed as a fixed-point
Hamiltonian in the same way that the toric code is a fixed-point Hamiltonian of
\(\ZZ_2\) topological order.  Perturbing the Hamiltonian generally destroys
exact solvability, but small local perturbations cannot immediately remove the
topological phase as long as the bulk gap remains open.

\section{Summary and comparison with the toric code}
\label{sec:quantum-double-summary}

The quantum double construction replaces every toric-code \(\ZZ_2\) ingredient
by its finite-group analogue.

{\small
\begin{longtable}{L{0.27\textwidth}L{0.30\textwidth}L{0.32\textwidth}}
\caption{Toric code versus finite-group quantum double models.}
\label{tab:toric-code-vs-quantum-double}\\
\toprule
Structure & Toric code & Quantum double \(D(G)\) \\
\midrule
\endfirsthead
\toprule
Structure & Toric code & Quantum double \(D(G)\) \\
\midrule
\endhead
Local edge Hilbert space & Qubit \(\CC[\ZZ_2]\) & Group algebra \(\CC[G]\) \\
Vertex term & \(\ZZ_2\) star constraint & Gauge-invariance projector \(A_v=\abs{G}^{-1}\sum_h A_v^h\) \\
Plaquette term & \(\ZZ_2\) zero-flux constraint & Flatness projector \(B_p^e\) \\
Ground states & Flat \(\ZZ_2\) gauge fields modulo gauge transformations & Flat \(G\)-connections modulo gauge transformations \\
Anyon labels & \(1,e,m,\epsilon\) & Pairs \((C,\rho)\) with \(C\) a conjugacy class and \(\rho\in\operatorname{Irr}(Z_c)\) \\
Abelian or non-Abelian & Abelian & Abelian iff \(G\) is Abelian; generally non-Abelian \\
Total quantum dimension & \(2\) & \(\abs{G}\) \\
\bottomrule
\end{longtable}}

\boxedpar{%
The toric code is not an isolated construction.  It is the \(G=\ZZ_2\) member of
the finite-group quantum double family.  The general replacement is
\[
\ZZ_2\text{ gauge theory}
\quad\longrightarrow\quad
G\text{ lattice gauge theory},
\]
with electric charges, magnetic fluxes, and dyons unified by the representation
theory of \(D(G)\).
}

The next chapter studies a different kind of two-dimensional exact solvability:
the Kitaev honeycomb model.  There the Hamiltonian is not a commuting-projector
Hamiltonian.  Instead, spin interactions fractionalize into free Majorana
fermions moving in a static \(\ZZ_2\) gauge field.
